\section{Critical growth, source reserve, and chord density}
\label{r:sec:critical-reserve-density}

The source work at critical order is \(\ip{u}{\Lambda g}\), where
\(g=\Pp f\). Its spatial derivatives can all be placed on \(g\), so the
kinetic bound controls how much this term can still add before a prescribed
finite time. Keeping that remaining contribution with the critical height
gives a positive quantity whose growth is governed by the nonlinear
production of the velocity.

Let \(u\) be a classical history on \([0,T)\), with the force smooth through
\(T\) as in Section~\ref{r:sec:introduction}. All calculations first take place
on strict classical intervals. Write
\begin{equation}
 \begin{aligned}
 \mathsf K_T&=\norm{u_0}_2+\int_0^T\norm{g(t)}_2\dd t,\\
 H(t)&=\tfrac12\norm{\Lambda^{1/2}u(t)}_2^2,
 &D(t)&=\norm{\Lambda^{3/2}u(t)}_2^2.
 \end{aligned}
 \label{r:crd-eq:height}
\end{equation}
The kinetic work satisfies
\[
 \int_0^t\norm{g(s)}_2\norm{u(s)}_2\dd s
 \le\int_0^t\mathsf K_s\mathsf K_s'\dd s
 =\tfrac12\bigl(\mathsf K_t^2-\norm{u_0}_2^2\bigr).
\]
The norm inequality and the energy identity give
\begin{equation}
 \norm{u(t)}_2\le\mathsf K_t,\qquad
 \norm{u(t)}_2^2+2\nu\int_0^t\norm{\nabla u(s)}_2^2\dd s
 \le\mathsf K_t^2\le\mathsf K_T^2.
 \label{r:crd-eq:kinetic}
\end{equation}
At critical order, put \(P_H=\mathcal P_{0,1/2}\) and
\(W=\operatorname{Re}\ip{u}{\Lambda g}\). The exact balance is
\[
 H'+\nu D=P_H+W,\qquad
 |W(t)|\le b_f(t):=\mathsf K_T\norm{\Lambda g(t)}_2.
\]
The source assumptions give \(b_f\in L^1(0,T)\). Its remaining integral
decreases at the rate \(b_f\), so define
\begin{equation}
 \mathscr R_f(t)=\int_t^T b_f(s)\dd s,\qquad
 \Phi=H+\mathscr R_f,\qquad Z_f=b_f-W.
 \label{r:crd-eq:reserve}
\end{equation}
Then
\begin{equation}
 \Phi'+\nu D+Z_f=P_H,\qquad
 Z_f\ge0,\qquad \Phi\ge H\ge0.
 \label{r:crd-eq:reserve-balance}
\end{equation}
The known function \(\mathscr R_f\) is absolutely continuous and tends to
zero at \(T\). Neither statement requires a terminal value of \(H\).
The production in \eqref{r:crd-eq:reserve-balance} is still evaluated on the
forced velocity. The force and its pressure contribution remain in the
momentum equation.

\subsection{Where critical growth is distributed}

The critical norm measures velocity differences over every separation.
For a center \(z\), a chord \(y\), and \(0\le\theta\le1\), put
\[
 x_+=z+(1-\theta)y,\qquad x_-=z-\theta y,\qquad
 \delta u_\theta=u(x_+)-u(x_-),\qquad \widehat y=y/|y|.
\]
The positive tensor
\begin{equation}
 \mathsf T_u(z)=\frac1{\pi^2}\int_0^1\int_{\R^3}
 \frac{|\delta u_\theta(z,y)|^2}{|y|^4}
 \widehat y\otimes\widehat y\dd y\dd\theta,\qquad
 \tau=\tr\mathsf T_u
 \label{r:crd-eq:tensor}
\end{equation}
weights each chord orientation by the squared magnitude of its velocity difference.
The fractional Fourier identity and the complete transport pairing give
\begin{equation}
 \int_{\R^3}\tau\dd z=2\norm{\Lambda^{1/2}u}_2^2=4H,\qquad
 P_H=-\int_{\R^3}S:\mathsf T_u\dd z,
 \quad S=\tfrac12(\nabla u+\nabla u^T).
 \label{r:crd-eq:current}
\end{equation}
The transport identity follows by integration in the chord variable and
averaging the strain along that chord; the local derivation below retains
the resulting redistribution flux.

Where \(\tau>0\), define \(r=-S:\mathsf T_u/\tau\), and set \(r=0\)
where \(\tau=0\). Positivity of \(\mathsf T_u\) gives
\[
 |r|\le\norm S_{\mathrm{op}}\le |S|_F,\qquad
 P_H=\int\tau r\dd z.
\]
Thus the production samples strain with the nonlocal density \(\tau\).
For \(H>0\), use the probability measure and its positive sampled rate
\[
 \dd\mu_t=\frac{\tau(z,t)\dd z}{4H(t)},\qquad
 \overline r_+(t)=\int r_+\dd\mu_t.
\]
Set \(\overline r_+=0\) when \(H=0\). Finite energy then forces \(u=0\),
so \(\tau=P_H=D=W=0\); no probability measure is assigned to zero mass.
Equation~\eqref{r:crd-eq:reserve-balance} gives
\begin{equation}
 \Phi'\le P_H\le\int\tau r_+\dd z
 =4H\overline r_+\le4\Phi\overline r_+.
 \label{r:crd-eq:sampled-growth}
\end{equation}
On an interval \([a,b]\) where \(\Phi>0\), a doubling requires
\[
 4\int_a^b\overline r_+(t)\dd t
 \ge\log\frac{\Phi(b)}{\Phi(a)}\ge\log2.
\]
The rate itself has an ordinary-volume square bound. For each eigenvalue
\(\lambda_i\) of the trace-free strain \(S\), the other two sum to
\(-\lambda_i\), so
\[
 |S|_F^2\ge\lambda_i^2+\tfrac12\lambda_i^2
 \quad\text{for every }i,\qquad
 |r|^2\le\norm S_{\mathrm{op}}^2\le\tfrac23|S|_F^2.
\]
Incompressibility gives \(\norm S_2^2=\norm{\nabla u}_2^2/2\), hence
\begin{equation}
 \int_0^T\int_{\R^3}|r|^2\dd z\dd t
 \le\frac{\mathsf K_T^2}{6\nu}.
 \label{r:crd-eq:rate-budget}
\end{equation}
The integral through \(T\) is the monotone limit of strict-subinterval
integrals.

To compare this volume bound with the density that samples it, put
\begin{equation}
 \rho_f=\frac{\tau}{4\Phi},\qquad
 V_f^{-1}(t)=\int_{\{r>0\}}\rho_f(z,t)^2\dd z,\qquad
 G_f=\{t:\Phi'(t)>0\}.
 \label{r:crd-eq:volume}
\end{equation}
These definitions apply where \(\Phi>0\); set \(\rho_f=0\) when
\(\Phi=0\). Its mass is \(H/\Phi\le1\). If
\(V_{\mathrm{work}}^{-1}=\int_{\{r>0\}}(\tau/(4H))^2\dd z\) for
\(H>0\), then
\[
 V_f^{-1}=(H/\Phi)^2V_{\mathrm{work}}^{-1}.
\]
At \(H=0<\Phi\), equation~\eqref{r:crd-eq:reserve-balance} gives
\(\Phi'=-b_f\le0\), so such times add no positive growth.

On a positive interval, discard the negative part of the logarithmic
growth and apply Cauchy--Schwarz in space and then in time:
\begin{equation}
 \log\frac{\Phi(b)}{\Phi(a)}
 \le4
 \left(\int_{[a,b]\cap G_f}\int r_+^2\dd z\dd t\right)^{1/2}
 \left(\int_{[a,b]\cap G_f}V_f^{-1}\dd t\right)^{1/2}.
 \label{r:crd-eq:growth-volume}
\end{equation}
A finite inverse-volume integral on the growth set thus bounds \(\Phi\)
from any positive initial value, and bounds \(H\) with it.
For \(N\) disjoint positive doubling intervals \(I_j\), write
\[
 A_j=\int_{I_j\cap G_f}\int r_+^2\dd z\dd t,\qquad
 C_j=\int_{I_j\cap G_f}V_f^{-1}\dd t.
\]
Then \(A_jC_j\ge(\log2)^2/16\). If \(\mathsf K_T>0\), a second
Cauchy--Schwarz inequality and \eqref{r:crd-eq:rate-budget} give
\begin{equation}
 \sum_{j=1}^N C_j
 \ge\frac{(\log2)^2N^2}{16\sum_{j=1}^N A_j}
 \ge\frac{3\nu(\log2)^2}{8\mathsf K_T^2}N^2.
 \label{r:crd-eq:repeated-growth}
\end{equation}
Infinite actions are read in the extended sense. If \(\mathsf K_T=0\),
the velocity is zero and no positive doubling occurs. Repeated growth
requires increasing concentration of the density that samples strain.
The evolution of that density must retain the endpoints of each chord.

\subsection{The local density equation and its source}

Write \(p=p_u+p_f\), with
\[
 -\Delta p_u=\partial_i u_j\,\partial_j u_i,\qquad
 \nabla p_f=(I-\Pp)f.
\]
Combining \(p_f\) with the prescribed force gives
\(D_tu=\nu\Delta u-\nabla p_u+g\). Define
\[
 a_\pm=u(x_\pm)-u(z),\qquad
 D_z=\partial_t+u(z)\cdot\nabla_z,
\]
and use the chord integral
\[
 \mathcal K[F](z)=\frac1{\pi^2}\int_0^1\int_{\R^3}
             \frac{F(z,\theta,y)}{|y|^4}\dd y\dd\theta.
\]
A fixed chord translated by the center velocity does not move with either
endpoint. Substituting the momentum equation at both endpoints gives
\begin{equation}
 \begin{aligned}
 D_z\delta u_\theta&=\nu\Delta_z\delta u_\theta
                  +\mathcal R+\delta g_\theta,\\
 \mathcal R&=-\delta\nabla p_u
 -(a_+\cdot\nabla)u(x_+)+(a_-\cdot\nabla)u(x_-).
 \end{aligned}
 \label{r:crd-eq:chord-momentum}
\end{equation}
Since \(\tau=\mathcal K[|\delta u_\theta|^2]\), the product rule yields
\begin{equation}
 \begin{aligned}
 (D_z-\nu\Delta_z)\tau
   &=\mathcal I_\tau-2\nu Q_\tau+\mathcal F_\tau,\\
 \mathcal I_\tau&=2\mathcal K[\delta u_\theta\cdot\mathcal R],\\
 Q_\tau&=\mathcal K[|\delta\nabla u_\theta|^2],\qquad
 \mathcal F_\tau=2\mathcal K[\delta u_\theta\cdot\delta g_\theta].
 \end{aligned}
 \label{r:crd-eq:density}
\end{equation}
Pressure and endpoint transport remain together in \(\mathcal I_\tau\).
The direct force term has an absolute spatial bound. Cauchy--Schwarz
over all centers, chords, and center parameters gives
\begin{equation}
 \norm{\mathcal F_\tau(t)}_1
 \le4\norm{\Lambda^{1/2}u(t)}_2\norm{\Lambda^{1/2}g(t)}_2
 =4\sqrt{2H(t)}\,\norm{\Lambda^{1/2}g(t)}_2.
 \label{r:crd-eq:local-source}
\end{equation}

The time integral can be bounded before a supremum of \(H\) is known.
For \(0<s\le1\), Fourier interpolation gives
\[
 E_s^{1/s}\le
 2^{-1/s}\norm u_2^{2(1-s)/s}\norm{\nabla u}_2^2,
 \qquad E_s=\tfrac12\norm{\Lambda^su}_2^2.
\]
Using \eqref{r:crd-eq:kinetic},
\begin{equation}
 \int_0^T E_s^{1/s}\dd t
 \le2^{-1/s-1}\nu^{-1}\mathsf K_T^{2/s},\qquad
 \int_0^T H^2\dd t\le\frac{\mathsf K_T^4}{8\nu}.
 \label{r:crd-eq:low-height}
\end{equation}
When \(\mathsf K_T=0\), both sides vanish. Applying H\"older to
\eqref{r:crd-eq:local-source} now gives
\begin{equation}
 \int_0^T\norm{\mathcal F_\tau(t)}_1\dd t
 \le2^{7/4}\mathsf K_T\nu^{-1/4}
       \norm{\Lambda^{1/2}g}_{L^{4/3}_tL^2_x}.
 \label{r:crd-eq:local-source-budget}
\end{equation}
The manuscript's smooth source class supplies this norm as well as the
\(L^1_tH^1_x\) norm used in \eqref{r:crd-eq:reserve}.
The direct chord forcing has finite absolute mass through \(T\), without
a terminal critical-height bound.

The integrated density equation has the normalization
\begin{equation}
 \int Q_\tau\dd z=2D,\qquad
 \int\mathcal F_\tau\dd z=4W,\qquad
 \int\mathcal I_\tau\dd z=4P_H.
 \label{r:crd-eq:integrated-density}
\end{equation}
For each fixed chord, \(\delta u_\theta\) is divergence free in \(z\),
so its global pairing with \(\delta\nabla p_u\) vanishes.
Integrating \eqref{r:crd-eq:density} gives
\(4H'=4P_H-4\nu D+4W\).
Pressure can still redistribute density locally. To expose that
redistribution, its flux must be retained before spatial integration.

\subsection{Intrinsic production and redistribution}

Endpoint advection acts on the chord coordinates through
\[
 c_\theta\cdot\nabla_z+\delta u_\theta\cdot\nabla_y,\qquad
 c_\theta=\theta u(x_+)+(1-\theta)u(x_-).
\]
Both coordinate divergences vanish by incompressibility. The intrinsic
integrand in \eqref{r:crd-eq:density} can be written
\[
 \nabla_z\cdot\bigl[(u(z)-c_\theta)|\delta u_\theta|^2
               -2\delta p_\theta\,\delta u_\theta\bigr]
 -\nabla_y\cdot\bigl(\delta u_\theta|\delta u_\theta|^2\bigr),
\]
where \(\delta p_\theta=p_u(x_+)-p_u(x_-)\).
The first term already has a flux in the center coordinate:
\begin{equation}
 J_0=\mathcal K\bigl[(u(z)-c_\theta)|\delta u_\theta|^2
               -2\delta p_\theta\,\delta u_\theta\bigr].
 \label{r:crd-eq:first-flux}
\end{equation}
Integrating the second term against the radial kernel produces
\(-4|\delta u_\theta|^2(\delta u_\theta\cdot\widehat y)/|y|\).
The velocity difference is the integral of its gradient along the chord.
With \(s(z,y)=\widehat y\cdot S(z)\widehat y\),
\[
 \frac{\delta u_\theta\cdot\widehat y}{|y|}
 =\int_0^1s(z+(\alpha-\theta)y,y)\dd\alpha.
\]
Changing the center parameter preserves the endpoints:
\[
 \delta u_\alpha(z+(\alpha-\theta)y,y)
 =\delta u_\theta(z,y).
\]
Put \(F_\alpha(z,y)=|\delta u_\alpha(z,y)|^2s(z,y)\).
Its translation away from \(z\) has the exact divergence form
\[
 F_\alpha(z+h,y)-F_\alpha(z,y)
 =\nabla_z\cdot
 \left[h\int_0^1F_\alpha(z+\lambda h,y)\dd\lambda\right].
\]
With \(h_{\alpha\theta}=(\alpha-\theta)y\), define
\begin{equation}
 \begin{aligned}
 L(z)={}&\frac1{\pi^2}\int_{\R^3}\frac{\dd y}{|y|^4}
 \int_0^1\int_0^1 h_{\alpha\theta}
 \left[\int_0^1F_\alpha(z+\lambda h_{\alpha\theta},y)
                         \dd\lambda\right]\dd\alpha\dd\theta,\\
 J_{\mathrm{int}}={}&J_0-4L.
 \end{aligned}
 \label{r:crd-eq:full-flux}
\end{equation}
The untranslated part is \(S:\mathsf T_u\), giving
\begin{equation}
 \mathcal I_\tau=-4S:\mathsf T_u+\nabla_z\cdot J_{\mathrm{int}}
 =4\tau r+\nabla_z\cdot J_{\mathrm{int}}.
 \label{r:crd-eq:intrinsic-split}
\end{equation}
The complete local equation is
\begin{equation}
 \partial_t\tau+\nabla_z\cdot(u\tau-J_{\mathrm{int}})
 -\nu\Delta_z\tau
 =4\tau r-2\nu Q_\tau+\mathcal F_\tau.
 \label{r:crd-eq:local-conservation}
\end{equation}
The strain at the center controls the displayed production, while
\(J_{\mathrm{int}}\) retains pressure and the change between chord centers.

These operations can first be performed with
\(\varepsilon<|y|<R\). The boundary remainder in the chord integration
obeys
\[
 \norm{B_{\varepsilon,R}}_1
 \le C\bigl(\varepsilon\norm{\nabla u}_3^3
                 +R^{-2}\norm u_3^3\bigr)\longrightarrow0.
\]
No endpoint in \(\theta\) or \(\alpha\) is differentiated.
On strict smooth time windows the paired differences also give two
powers of \(|y|\) near zero, and the radial kernel is integrable at
infinity. These bounds justify the local identities in spatial
distributions.

The flux itself is absolutely integrable in space on such a window.
Write \(U=\norm u_2\), \(A=\norm{\nabla u}_2\), and \(B=\norm u_3\).
Fractional Sobolev embedding and the Fourier kernel identity yield
\[
 \int_{\R^3}\frac{\norm{\delta_yu}_3^2}{|y|^4}\dd y
 \le C\int_{\R^3}
 \frac{\norm{\Lambda^{1/2}\delta_yu}_2^2}{|y|^4}\dd y
 \le CA^2.
\]
Since \(\norm{u-c_\theta}_3\le2B\), the transport part of \(J_0\)
has \(L^1\) norm at most \(CBA^2\).

For \(L\), split the chords at length \(\ell>0\).
The estimate
\(\norm{\delta_yu}_4^2\le C|y|^{1/2}A^2\) gives
\(C\ell^{1/2}\norm S_2A^2\) for short chords.
For long chords, average the squared endpoint velocities in the center
parameter. For an endpoint fraction \(0<a<1\), changing variables to
\(h=ay\) gives the kernel
\[
 k_\ell(h)=|h|^{-3}\min(1,|h|/\ell),\qquad
 \norm{k_\ell}_{6/5}^{6/5}=\frac{40\pi}{3}\ell^{-3/5}.
\]
H\"older and Young convolution inequalities bound this part by
\(C\ell^{-1/2}\norm S_2B^2\).
Choosing \(\ell=(B/A)^2\) for nonzero \(u\) gives
\(\norm L_1\le C\norm S_2AB\le CBA^2\).
For the zero field there is no division.

The pressure part of \(J_0\) satisfies
\[
 \norm{J_{0,\mathrm{pressure}}}_1
 \le4\norm{\Lambda^{1/2}p_u}_2\norm{\Lambda^{1/2}u}_2.
\]
The pressure multiplier is bounded, and Fourier Cauchy--Schwarz for
each product \(u_i u_j\) uses
\[
 \int_{\R^3}\frac{\dd\eta}{|\eta|^2|\xi-\eta|^2}
 =C|\xi|^{-1}.
\]
Multiplying by \(|\xi|\), integrating in \(\xi\), and applying Tonelli
gives \(\norm{\Lambda^{1/2}p_u}_2\le CA^2\).
Finally \(B\le CU^{1/2}A^{1/2}\) and
\(\norm{\Lambda^{1/2}u}_2\le U^{1/2}A^{1/2}\), so
\begin{equation}
 \norm{J_{\mathrm{int}}}_1\le CU^{1/2}A^{5/2},\qquad
 \int_0^T\norm{J_{\mathrm{int}}}_1^{4/5}\dd t
 \le C\nu^{-1}\mathsf K_T^{12/5}.
 \label{r:crd-eq:flux-budget}
\end{equation}
The time exponent is \(4/5\). A terminal \(L^1_tL^1_z\) flux estimate
does not follow from this bound.

\subsection{Diffusion of the normalized density}

Diffusion gives a local constraint on how sharply the density varies.
Cauchy--Schwarz inside the chord integral gives
\[
 |\nabla\tau|^2\le4\tau Q_\tau,\qquad
 -2\nu Q_\tau\le-\frac\nu2\frac{|\nabla\tau|^2}{\tau}
 \quad\hbox{where }\tau>0.
\]
The intrinsic production in \eqref{r:crd-eq:density} remains alongside
this favorable term. Its time integral cannot be discarded when estimating
the diffusion of \(\rho_f\).

On an interval with \(\Phi>0\), put
\[
 \psi=\sqrt{\rho_f},\qquad
 m=\int\rho_f\dd z=\frac H\Phi\le1,\qquad
 \mathcal J=\int|\nabla\psi|^2\dd z.
\]
Differentiating the normalization gives
\begin{equation}
 (D_z-\nu\Delta_z)\rho_f
 =\frac{\mathcal I_\tau+\mathcal F_\tau}{4\Phi}
 -\frac{\nu Q_\tau}{2\Phi}
 -\frac{\Phi'}{\Phi}\rho_f.
 \label{r:crd-eq:normalized-density}
\end{equation}
Where \(\tau>0\), its square root satisfies
\begin{equation}
 \begin{aligned}
 (D_z-\nu\Delta_z)\psi
 ={}&\frac{\mathcal I_\tau+\mathcal F_\tau}{8\Phi\psi}
 -\frac{\nu}{\psi}
     \left(\frac{Q_\tau}{4\Phi}-|\nabla\psi|^2\right)
 -\frac{\Phi'}{2\Phi}\psi.
 \end{aligned}
 \label{r:crd-eq:square-root}
\end{equation}
The bracket is nonnegative. The derivative of \(\Phi\) retains its full
sign, and pressure and endpoint transport remain inside \(\mathcal I_\tau\).

The norm derivative inequality extends across zeros by Sobolev approximation,
and \eqref{r:crd-eq:integrated-density} gives
\(\mathcal J\le D/(2\Phi)\).
Let \(S_6\) be a constant in
\(\norm h_6\le S_6\norm{\nabla h}_2\). Interpolation gives
\begin{equation}
 V_f^{-1}\le\norm\psi_4^4
 \le\min\left\{
 S_6^3m^{1/2}\mathcal J^{3/2},
 \frac{S_6^2\norm\tau_{3/2}\mathcal J}{4\Phi}
 \right\}.
 \label{r:crd-eq:density-interpolation}
\end{equation}
The first estimate uses the \(L^2\) and \(L^6\) norms of \(\psi\);
the second uses its \(L^3\) and \(L^6\) norms.
For the chord density,
\(\norm\tau_{3/2}\le C_\tau\norm{\nabla u}_2^2\):
Minkowski bounds this norm by
\(C\int |y|^{-4}\norm{\delta_yu}_3^2\dd y\), which is controlled by
the fractional estimate used for \(J_0\).
Substituting these bounds gives
\begin{equation}
 V_f^{-1}\le\min\left\{
 \frac{S_6^3H^{1/2}D^{3/2}}{2^{3/2}\Phi^2},
 \frac{C_\tau S_6^2\norm{\nabla u}_2^2D}{8\Phi^2}
 \right\}.
 \label{r:crd-eq:inverse-volume}
\end{equation}
Neither time integral on the right is bounded by
\eqref{r:crd-eq:kinetic} and \eqref{r:crd-eq:local-source-budget}.
Even a finite integral of \(\mathcal J\) would give only
\[
 \int(V_f^{-1})^{2/3}\dd t\le S_6^2\int\mathcal J\dd t,
\]
which has a smaller time exponent than the action in
\eqref{r:crd-eq:growth-volume}.

For this route, a sufficient additional estimate is
\[
 \int_{G_f}V_f^{-1}(t)\dd t<\infty
\]
with a bound derived from \(u_0,\nu,f\). The source reserve and the absolute
chord-source estimate control the direct external contribution.
The required intrinsic estimate must control how strain samples the density
under the complete production and redistribution in
\eqref{r:crd-eq:local-conservation}. The rate bound
\eqref{r:crd-eq:rate-budget}, the flux bound \eqref{r:crd-eq:flux-budget},
and the instantaneous diffusion estimates above do not supply that
time-integrated control.
